\documentclass[11pt,twocolumn]{article}

\usepackage[a4paper,margin=0.75in,columnsep=0.3in]{geometry}
\usepackage[T1]{fontenc}
\usepackage{amsmath,amssymb,amsthm,mathtools}
\usepackage{bm}
\usepackage{graphicx}
\usepackage{booktabs}
\usepackage{array}
\usepackage{enumitem}
\usepackage{float}
\usepackage[numbers,sort&compress]{natbib}
\usepackage{hyperref}
\hypersetup{colorlinks=true,linkcolor=blue,citecolor=blue,urlcolor=blue}

\newtheorem{theorem}{Theorem}
\newtheorem{definition}{Definition}
\newtheorem{postulate}{Postulate}
\newtheorem{corollary}{Corollary}
\newtheorem{proposition}{Proposition}

\title{Unified Dynamics III: Dimensional Growth, Spectral Geometry, and the Mass Spectrum from the Tetrahedral Closure Ladder}
\author{%
Benjamin Walker Mayes\thanks{*Contact author: ben.w.mayes@gmail.com}\\[4pt]
{\small Unaffiliated, Cedar Park, Texas 78613, USA}\\[8pt]
Arthur Eugene Mayes\\
{\small Loyola University School of Law, USA}\\[4pt]
Dustin Huff\\
{\small Unaffiliated, USA}\\[4pt]
Brayden Ross Sanders\\
{\small Unaffiliated, USA}\\[4pt]
Jeffrey A.\ Dewey\\
{\small Unaffiliated, USA}\\[4pt]
Brandon Bevin\\
{\small Georgia Southern University, USA}\\[4pt]
A.\ Christopher Beckingham\\
{\small Queen's University, Canada}\\[4pt]
Jonathan Engle\\
{\small Unaffiliated, USA}%
}
\date{}

\begin{document}
\maketitle

\begin{abstract}
Building on the tetrahedral transport--closure framework of Unified Dynamics (UD), we prove that $d=3$ is the unique stable minimal closure dimension among all simplex dimensions, derive the spectral dimension $d_s=4/3$ of the Sierpi\'nski tetrahedron from decimation, and extend the mass spectrum with new predictions.
The $d$-simplex $K_{d+1}$ has state-space dimension $N(d)=(d{+}1)!\cdot 2^d$, overlap deficit $\Delta(d)=d(d{+}1)/2{+}1$, and master scalar $x(d)=\Delta/N$.
We show that $d=3$ uniquely satisfies four simultaneous conditions: face--vertex balance, perturbative stability, binding viability, and optimal coherence retention.
For $d\geq 4$, the nonlinear binding coupling falls below the critical threshold and lone primitives cannot self-bind; stable higher-dimensional structures require composites of $d=3$ closures with enhanced coupling $\lambda_{\mathrm{NL}}^{\mathrm{comp}}=n_c^2\lambda_{\mathrm{NL}}(d)$.
The coherence ceiling at $n_c^{\max}=61$ caps the dimensional ladder at $d_{\max}\approx 7$.
Additionally, we derive the cross-coupling $g_{\chi\sigma}=\Delta(\Delta{-}1)/(2^{E+1}\pi)$ from $K_4$ combinatorics, prove confinement via the string tension $\sigma=9/(4\pi^2 a_0)$, predict the dark-matter relic fraction $\Omega_{\mathrm{DM}}=43/228\approx 19\%$, present neutrino masses $m_{\nu_1}\approx 35$~meV and PMNS angles from generation-spacing mismatch, and sharpen the muon mass prediction to $0.2\%$.
\end{abstract}

%═══════════════════════════════════════════════════════════
\section{Introduction}
\label{sec:intro}
%═══════════════════════════════════════════════════════════

A foundational question for any theory of emergent geometry is: why three spatial dimensions?
The companion papers established a transport--closure framework on the tetrahedron $K_4$~\cite{mayes2026ud} and derived its cosmological consequences~\cite{mayes2026cosmo}.
The present paper addresses the dimensional question directly by generalizing the $K_4$ construction to the $d$-simplex $K_{d+1}$ and proving that $d=3$ is the unique stable closure dimension.

We also close several gaps from the earlier papers: the spectral dimension of the IFS lattice is derived exactly ($d_s=4/3$), the cross-coupling $g_{\chi\sigma}$ is derived rather than stated, confinement is proved from the mass gap, the dark-matter abundance receives a structural prediction, and the matter sector is extended to neutrinos and sharpened for the muon.

%═══════════════════════════════════════════════════════════
\section{Simplex Combinatorics in Arbitrary Dimension}
\label{sec:simplex}
%═══════════════════════════════════════════════════════════

The $d$-dimensional simplex $K_{d+1}$ has $V=d{+}1$ vertices, $E=d(d{+}1)/2$ edges, $F_2=\binom{d+1}{3}$ triangular faces, and symmetry group $S_{d+1}$.
The UD state-space dimension and overlap deficit generalize to~\cite{hutchinson1981}
\begin{align}
N(d) &= (d{+}1)!\cdot 2^d, \label{eq:Nd}\\
\Delta(d) &= \tbinom{d{+}1}{2}+1 = \tfrac{d(d{+}1)}{2}+1, \label{eq:Dd}\\
x(d) &= \Delta(d)/N(d). \label{eq:xd}
\end{align}
The coherence retention is $\eta_R(d)=1-x(d)$, the binding coupling is $\lambda_{\mathrm{NL}}(d)=[\Delta(\Delta{-}1)]^2/(2^{2E+2}\pi^2 J_\sigma)$, and the critical binding threshold is $\lambda_{\mathrm{crit}}(d)=d/N(d)$, where the factor $d$ accounts for the increased kinetic energy cost of self-localization in $d$ dimensions.

%═══════════════════════════════════════════════════════════
\section{$D=3$ Uniqueness}
\label{sec:d3}
%═══════════════════════════════════════════════════════════

\begin{theorem}[Three-dimensional uniqueness]
\label{thm:d3}
$d=3$ is the unique dimension satisfying all four conditions:
\begin{enumerate}[label=\textnormal{(\roman*)},leftmargin=2em,itemsep=2pt]
\item Face--vertex balance: $F_2=V$, requiring $\binom{d+1}{3}=d+1$, solved uniquely by $d=3$.
\item Perturbative stability: $x(d)\in(0.01,\,0.1)$.
\item Binding viability: $\lambda_{\mathrm{NL}}(d)>\lambda_{\mathrm{crit}}(d)=d/N(d)$.
\item Optimal coherence: $\eta_R(d)\in(0.9,\,0.99)$.
\end{enumerate}
\end{theorem}

\textit{Proof sketch.}
Condition~(i): $\binom{d+1}{3}=d+1$ gives $d^2-d-6=0$, so $d=3$.
Condition~(ii): $x(1)=1/2$, $x(2)=1/6$, $x(3)=7/192\approx 0.036$, $x(4)\approx 0.006$; only $d=3$ lies in $(0.01,0.1)$.
Condition~(iii): $\lambda_{\mathrm{NL}}(3)=0.022>3/192=0.0156$; $\lambda_{\mathrm{NL}}(4)=5.9\times 10^{-4}<4/1920=2.1\times 10^{-3}$. Binding holds for $d\leq 3$, fails for $d\geq 4$.
Condition~(iv): $\eta_R(3)=185/192\approx 0.964$; $\eta_R(2)=5/6\approx 0.83$ (too lossy); $\eta_R(4)\approx 0.994$ (insufficient selection).

%═══════════════════════════════════════════════════════════
\section{$D=4$ Burden Regime}
\label{sec:d4}
%═══════════════════════════════════════════════════════════

\begin{theorem}[$D=4$ burden]
\label{thm:d4}
The pentachoral regime $d=4$ has $N=1920$, $\Delta=11$, $E=10$, $\lambda_{\mathrm{NL}}=5.9\times 10^{-4}<\lambda_{\mathrm{crit}}=4/1920=2.1\times 10^{-3}$.
Lone $d=4$ primitives cannot form bound states.
\end{theorem}

\textit{Proof sketch.}
Transport burden $B(d)=d(d{+}1)/4$ gives $B(4)=5$ vs $B(3)=3$ (67\% increase).
The symmetry group $S_5$ ($|S_5|=120$) introduces sorting, coordination, expression, and renewal-deficit burdens.
With $x(4)=0.006$, selection pressure is $6\times$ weaker than $d=3$, preventing evolutionary self-correction for lone primitives.

%═══════════════════════════════════════════════════════════
\section{$D\geq 5$: Composite-First Stabilization}
\label{sec:d5}
%═══════════════════════════════════════════════════════════

\begin{theorem}[Composite binding enhancement]
\label{thm:composite}
A composite of $n_c$ tetrahedral ($d=3$) closures embedded in $d\geq 4$ space has enhanced coupling $\lambda_{\mathrm{NL}}^{\mathrm{comp}}=n_c^2\lambda_{\mathrm{NL}}(d)$ and effective selection pressure $x_{\mathrm{eff}}=n_c\,x(d)$.
For $d=4$, binding requires $n_c\geq 2$; strong stability requires $n_c\geq 8$.
\end{theorem}

Three composite classes emerge:

\textbf{Seed} ($n_c=4$): $\lambda^{\mathrm{comp}}/\lambda_{\mathrm{crit}}=4.5$, $x_{\mathrm{eff}}=0.023$.
First admissible opener; high residual bonding hunger.

\textbf{Stable composite} ($n_c=8$): ratio $=18$, $x_{\mathrm{eff}}=0.046$.
Shell-satisfied; low bonding hunger; first reusable unit.

\textbf{Shell molecule} ($n_c\geq 16$): ratio $\geq 70$, $x_{\mathrm{eff}}\geq 0.09$.
Bonded composites with $B_{\mathrm{molecule}}<2B_{\mathrm{composite}}$; first molecular-scale higher-dimensional regime.

The single lift mechanism governs each stage: rigid kernel $\to$ $\Omega$-tether $\to$ open channel $\to$ field $\to$ face completion $\to$ gate $\to$ next admissible geometry.

%═══════════════════════════════════════════════════════════
\section{Coherence Ceiling}
\label{sec:ceiling}
%═══════════════════════════════════════════════════════════

\begin{theorem}[Dimensional ceiling]
\label{thm:ceiling}
The maximum composite size is $n_c^{\max}=\lfloor N(3)\eta_R(3)/B(3)\rfloor=\lfloor 185/3\rfloor=61$.
The dimensional ladder terminates at $d_{\max}\approx 7$, where $n_c^{\min}(d)$ exceeds $61$.
\end{theorem}

\textit{Proof sketch.}
Each $d=3$ closure in the composite consumes $B(3)=3$ carrier units.
The total carrier budget is $N(3)\eta_R(3)=185$.
Maximum composites: $185/3\approx 61$.
Since $n_c^{\min}(d)$ grows as $\sim\sqrt{(d+1)!}$, the bound $n_c^{\min}\leq 61$ is saturated near $d=7$.
Beyond $d_{\max}$, composites rupture into $d=3$ remnants, incoherent residue, and kernel-bearing shards.

%═══════════════════════════════════════════════════════════
\section{Spectral Dimension}
\label{sec:spectral}
%═══════════════════════════════════════════════════════════

\begin{theorem}[Spectral dimension]
\label{thm:ds}
The spectral dimension of the Sierpi\'nski tetrahedron is $d_s=4/3$.
\end{theorem}

\textit{Proof sketch.}
By the decimation method~\cite{rammal1984}, the resistance scaling of the $d=3$ Sierpi\'nski gasket is $\rho=2$.
Then $d_s=2\ln V/\ln(V\rho)=2\ln 4/\ln 8=4/3$.

The IFS loop-integral correction factor for renormalization group calculations is $F_{\mathrm{IFS}}=d_s/4=1/3$.
This modifies the standard beta-function coefficients: $b_i^{\mathrm{IFS}}=b_i^{\mathrm{SM}}\cdot F_{\mathrm{IFS}}\cdot(\text{threshold corrections})$.

%═══════════════════════════════════════════════════════════
\section{Cross-Coupling Derivation}
\label{sec:coupling}
%═══════════════════════════════════════════════════════════

\begin{theorem}[Cross-coupling from $K_4$]
\label{thm:gxs}
The closure--framing cross-coupling is
\begin{equation}
g_{\chi\sigma} = \frac{\Delta(\Delta{-}1)}{2^{E+1}\pi} = \frac{21}{64\pi},
\label{eq:gxs}
\end{equation}
giving $\lambda_{\mathrm{NL}}=g_{\chi\sigma}^2/J_\sigma = 441/(2048\pi^2)$.
\end{theorem}

\textit{Proof sketch.}
The numerator $\Delta(\Delta{-}1)/2=\binom{7}{2}=21$ counts pairwise deficit-mode couplings.
The denominator $2^E=64$ is the binary edge-state space.
The factor $\pi$ is the consumed-loop phase per face.

%═══════════════════════════════════════════════════════════
\section{Confinement from the Mass Gap}
\label{sec:confinement}
%═══════════════════════════════════════════════════════════

\begin{theorem}[Confinement]
\label{thm:confine}
The mass gap $\Delta_{\mathrm{mass}}\geq 3/\pi$ on the IFS lattice implies an area-law Wilson loop with string tension $\sigma=9/(4\pi^2 a_0)$.
\end{theorem}

\textit{Proof sketch.}
The mass gap eliminates massless gauge modes~\cite{wilson1974}.
Correlation functions decay as $\langle A_\mu(x)A_\nu(0)\rangle\sim e^{-\Delta_{\mathrm{mass}}|x|}$.
For a Wilson loop of area $A$, $W(C)=\langle\mathrm{Tr}\,U_C\rangle\sim e^{-\sigma A}$ with $\sigma=\kappa_C\Delta_{\mathrm{mass}}/a_0 = (3/(4\pi))(3/\pi)/a_0=9/(4\pi^2 a_0)$.

%═══════════════════════════════════════════════════════════
\section{Dark-Matter Abundance}
\label{sec:dm}
%═══════════════════════════════════════════════════════════

\begin{proposition}[Relic fraction]
\label{prop:dm}
The dark-matter fraction of total matter density is
\begin{equation}
\frac{\Omega_{\mathrm{DM}}}{\Omega_{\mathrm{total}}} = \frac{x\,S_0}{x\,S_0+(1{-}x)} = \frac{43}{228}\approx 0.189.
\end{equation}
\end{proposition}

\textit{Argument.}
Neutral packets have mass per site $\nu_0 S_0$ (threshold cost); charged packets have mass per site $\nu_0$.
The mass-weighted neutral fraction is $x S_0/(x S_0+1{-}x)=43/228$.
The measured value $\Omega_{\mathrm{DM}}/\Omega_{\mathrm{total}}\approx 0.27$ is $30\%$ above the structural prediction; the discrepancy may close with dressing corrections.

%═══════════════════════════════════════════════════════════
\section{Mass Spectrum Extensions}
\label{sec:mass}
%═══════════════════════════════════════════════════════════

\textbf{Muon mass refinement.}
Sub-leading IFS corrections improve the muon prediction~\cite{mayes2026ud}:
\begin{equation}
\frac{m_\mu}{m_e} = 4^4\!\left(\frac{185}{192}\right)^{\!6}\!\left(1+\frac{\Delta{+}C^*}{NE}+\frac{F}{NEd}\right)\approx 206.3,
\label{eq:muon}
\end{equation}
where $C^*=7/8$ and $F=4$.
The measured ratio $206.77$ is matched to $0.2\%$.

\textbf{Neutrino masses.}
Neutrinos couple at second order in $x$ (two overlap-deficit insertions for a framing-neutral particle) with consumed-loop suppression $1/\pi$:
\begin{equation}
m_{\nu_1} = \frac{\nu_0 x^2}{\pi}\approx 35\text{ meV}.
\label{eq:neutrino}
\end{equation}
This is consistent with cosmological bounds $\sum m_\nu<0.12$~eV~\cite{planck2018}.

\textbf{PMNS mixing angles.}
The mismatch between charged-lepton generation spacings ($12{:}6$ IFS levels, ratio $2{:}1$) and uniform neutrino spacings ($1{:}1$) gives:
\begin{align}
\theta_{12} &= \arctan(1/2)\approx 26.6^\circ &(\text{meas.\ }33.4^\circ),\label{eq:t12}\\
\theta_{23} &= 45^\circ &(\text{meas.\ }42\text{--}50^\circ),\label{eq:t23}\\
\theta_{13} &= \theta_{12}/4\approx 6.6^\circ &(\text{meas.\ }8.6^\circ).\label{eq:t13}
\end{align}

%═══════════════════════════════════════════════════════════
\section{IFS Level Hierarchy}
\label{sec:hierarchy}
%═══════════════════════════════════════════════════════════

The lattice spacing $a_0=\sqrt{\kappa_C}/M_{\mathrm{Pl}}$ places the IFS at the Planck scale.
The energy per site $\nu_0\approx 83$~keV sits at IFS level $\sim 102$ below the Planck cutoff.
The electroweak scale $v_{\mathrm{EW}}\approx 246$~GeV sits at level $\sim 85$.
The full hierarchy is:
level~0 (Planck) $\to$ level~${\sim}85$ (electroweak) $\to$ level~${\sim}102$ ($\nu_0$) $\to$ level~${\sim}130$ (QCD) $\to$ level~${\sim}190$ (Hubble).

%═══════════════════════════════════════════════════════════
\section{Predictions and Falsifiers}
\label{sec:predictions}
%═══════════════════════════════════════════════════════════

\begin{table}[t]
\centering
\caption{Predictions from the dimensional growth theorem.}
\label{tab:pred3}
\small
\begin{tabular}{@{}l@{\hskip 5pt}l@{\hskip 5pt}l@{\hskip 5pt}l@{}}
\toprule
Quantity & Prediction & Meas.\ & Acc.\ \\
\midrule
$d_{\mathrm{stable}}$ & $3$ & $3$ & exact \\
$d_s$ & $4/3$ & --- & struct.\ \\
$d_{\max}$ & $\approx 7$ & --- & testable \\
$\sigma$ & $9/(4\pi^2 a_0)$ & --- & struct.\ \\
$g_{\chi\sigma}$ & $21/(64\pi)$ & --- & struct.\ \\
$\Omega_{\mathrm{DM}}$ & $43/228$ & $0.27$ & 30\% \\
$m_{\nu_1}$ & $35$ meV & $<120$ meV & consist.\ \\
$\theta_{23}$ & $45^\circ$ & $42$--$50^\circ$ & consist.\ \\
$\theta_{13}$ & $6.6^\circ$ & $8.6^\circ$ & 23\% \\
$m_\mu/m_e$ & $206.3$ & $206.8$ & 0.2\% \\
\bottomrule
\end{tabular}
\end{table}

Falsifiers:
(F12)~a fourth generation of fermions discovered;
(F13)~stable lone-primitive $d\geq 4$ bound states observed;
(F14)~spectral dimension $d_s\neq 4/3$ measured on a tetrahedral fractal;
(F15)~$\Omega_{\mathrm{DM}}$ shown to be incompatible with $43/228$ after dressing corrections.

%═══════════════════════════════════════════════════════════
\section{Discussion}
\label{sec:discussion}
%═══════════════════════════════════════════════════════════

The central result of this paper---that $d=3$ is uniquely selected by four simultaneous conditions from simplex combinatorics---provides a structural answer to the dimensionality question.
The answer does not invoke anthropic reasoning~\cite{weinberg1987} or extra-dimensional compactification~\cite{polyakov1987}; it follows from the binding criterion $\lambda_{\mathrm{NL}}(d)>d/N(d)$, which is a property of the simplex alone.
The entropy bound~\cite{bekenstein1973} is automatically respected because the carrier budget $N\eta_R$ caps the information content at each level.

The dark-matter prediction $\Omega_{\mathrm{DM}}=43/228$ is a zero-parameter structural result using only locked constants.
The $30\%$ discrepancy with the measured $27\%$ may close with dressing corrections analogous to those that improved the muon mass from $0.9\%$ to $0.2\%$.

The spectral dimension $d_s=4/3$ from decimation~\cite{rammal1984} places the IFS loop integrals on exact footing.
Combined with threshold corrections from the five packet channels and the carrier excitation spectrum (requiring simulation), this framework is expected to resolve the remaining $5$--$8\times$ discrepancy in the strong coupling $\alpha_s$ at the $Z$ scale.

Open computations include the full RG flow with IFS beta functions, the $n_c^{\min}(d)$ series for $d=5{-}7$, and lattice thermodynamics for the relic abundance and baryogenesis washout.

%═══════════════════════════════════════════════════════════
\section{Conclusion}
\label{sec:conclusion}
%═══════════════════════════════════════════════════════════

The tetrahedral closure ladder determines the full dimensional ontology of Unified Dynamics: $d=3$ is uniquely stable, higher dimensions require composite promotion, and the coherence ceiling caps the ladder at $d\approx 7$.
The spectral dimension $d_s=4/3$, the cross-coupling $g_{\chi\sigma}=21/(64\pi)$, confinement with string tension $\sigma=9/(4\pi^2 a_0)$, the dark-matter fraction $43/228$, neutrino masses near $35$~meV, and the muon mass to $0.2\%$ are derived from $K_4$ combinatorics without free parameters.

\section*{Acknowledgments}
The authors thank iterative computational dialog and theory-development workflows for enabling the progressive consolidation of the Unified Dynamics framework.

\section*{Author contributions}
B.W.M.\ conceived the framework, derived all results, and wrote the manuscript.
A.E.M., D.H., B.R.S., J.A.D., B.B., A.C.B., and J.E.\ contributed through discussion, review, and validation.

\smallskip
\noindent\textbf{Data availability statement.}
This paper has no associated data.
This is a theoretical study and no experimental data were generated or analyzed.

\begin{thebibliography}{99}
\bibitem{mayes2026ud} B.~W.~Mayes et al., ``Unified Dynamics: Emergent geometry, gauge structure, and matter from a tetrahedral transport-closure substrate,'' Zenodo (2026), \href{https://doi.org/10.5281/zenodo.19582493}{doi:10.5281/zenodo.19582493}.
\bibitem{mayes2026cosmo} B.~W.~Mayes et al., ``Unified Dynamics II: Dark matter, baryogenesis, the cosmological constant, and black hole information from the tetrahedral framing phase transition,'' Zenodo (2026), \href{https://doi.org/10.5281/zenodo.19582493}{doi:10.5281/zenodo.19582493}.
\bibitem{rammal1984} R.~Rammal and G.~Toulouse, ``Random walks on fractal structures and percolation clusters,'' J.\ Physique Lett.\ \textbf{44}, L13--L22 (1983).
\bibitem{wilson1974} K.~G.~Wilson, ``Confinement of quarks,'' Phys.\ Rev.\ D \textbf{10}, 2445--2459 (1974).
\bibitem{planck2018} Planck Collaboration, ``Planck 2018 results.\ VI.\ Cosmological parameters,'' Astron.\ Astrophys.\ \textbf{641}, A6 (2020).
\bibitem{weinberg1987} S.~Weinberg, ``Anthropic bound on the cosmological constant,'' Phys.\ Rev.\ Lett.\ \textbf{59}, 2607--2610 (1987).
\bibitem{polyakov1987} A.~M.~Polyakov, \textit{Gauge Fields and Strings} (Harwood Academic, Chur, 1987).
\bibitem{hutchinson1981} J.~E.~Hutchinson, ``Fractals and self-similarity,'' Indiana Univ.\ Math.\ J.\ \textbf{30}, 713--747 (1981).
\bibitem{bekenstein1973} J.~D.~Bekenstein, ``Black holes and entropy,'' Phys.\ Rev.\ D \textbf{7}, 2333--2346 (1973).
\end{thebibliography}

\end{document}
